\documentclass[a4paper]{arrowhead}

\usepackage[yyyymmdd]{datetime}
\usepackage{etoolbox}
\usepackage[T1]{fontenc}
\usepackage[utf8]{inputenc}
\usepackage{multirow}
\usepackage{longtable}
\usepackage{amssymb}
\usepackage{float}
\usepackage{tikz}
\usetikzlibrary{positioning,arrows.meta,fit,calc,shapes.misc,shapes.geometric}

\renewcommand{\dateseparator}{-}
\setlength{\emergencystretch}{3em}
\lstset{basicstyle=\ttfamily\footnotesize,breaklines=true,columns=fullflexible,frame=single,rulecolor=\color{gray!50},backgroundcolor=\color{gray!6}}
\tikzset{
  ahsys/.style={draw=ArrowheadBlue,very thick,fill=ArrowheadBlue!6,rounded corners=2pt,minimum width=3.2cm,minimum height=1.1cm,align=center,font=\small\bfseries},
  ahext/.style={draw=ArrowheadGrey,thick,dashed,fill=gray!5,rounded corners=2pt,minimum width=3.2cm,minimum height=1.0cm,align=center,font=\small},
  ahsrv/.style={draw=ArrowheadPurple,thick,fill=ArrowheadPurple!8,rounded corners=6pt,minimum height=0.55cm,inner xsep=4pt,align=center,font=\scriptsize\ttfamily},
  ahcsm/.style={draw=ArrowheadTeal,thick,fill=ArrowheadTeal!8,rounded corners=6pt,minimum height=0.55cm,inner xsep=4pt,align=center,font=\scriptsize\ttfamily},
  ahdb/.style={draw=ArrowheadOrange,thick,fill=ArrowheadOrange!8,cylinder,shape border rotate=90,aspect=0.25,minimum height=0.9cm,align=center,font=\scriptsize},
  ahflow/.style={-{Stealth[length=2.2mm]},thick,ArrowheadGrey},
  ahlabel/.style={font=\scriptsize,fill=white,inner sep=1pt,align=center},
}

%% Special references
\newcommand{\mref}[1]{{\textcolor{ArrowheadPurple}{\hyperref[sec:model:#1]{#1}}}}
\newcommand{\oref}[1]{{\textcolor{ArrowheadBlue}{\hyperref[sec:operations:#1]{#1}}}}
\newcommand{\pdef}[1]{{\textcolor{ArrowheadGrey}{#1\label{sec:model:primitives:#1}\label{sec:model:primitives:#1s}\label{sec:model:primitives:#1es}}}}
\newcommand{\pref}[1]{{\textcolor{ArrowheadGrey}{\hyperref[sec:model:primitives:#1]{#1}}}}

\newrobustcmd\osubsection[3]{
  \addtocounter{subsection}{1}
  \addcontentsline{toc}{subsection}{\protect\numberline{\thesubsection}operation \textcolor{ArrowheadBlue}{#1}}
  \subsection*{
    \thesubsection\quad
    operation
    \textcolor{ArrowheadBlue}{#1}
    (\notblank{#2}{\mref{#2}}{})
    \notblank{#3}{: \mref{#3}}{}
  }
  \label{sec:operations:#1}
}
\newrobustcmd\msubsection[2]{
  \addtocounter{subsection}{1}
  \addcontentsline{toc}{subsection}{\protect\numberline{\thesubsection}#1 \textcolor{ArrowheadPurple}{#2}}
  \subsection*{\thesubsection\quad#1 \textcolor{ArrowheadPurple}{#2}}
  \label{sec:model:#2} \label{sec:model:#2s} \label{sec:model:#2es}
}

\begin{document}

%% Arrowhead Document Properties
\ArrowheadTitle{workflowLease}
\ArrowheadServiceID{workflowLease}
\ArrowheadType{Service Description}
\ArrowheadTypeShort{SD}
\ArrowheadVersion{5.4.0}
\ArrowheadDate{2026-11-23}
\ArrowheadAuthor{Documentation maintenance}
\ArrowheadStatus{REVIEW DRAFT}
\ArrowheadContact{Repository maintainers}
\ArrowheadFooter{\href{www.arrowhead.eu}{www.arrowhead.eu}}
\ArrowheadSetup
%%

%% Front Page
\begin{center}
  \vspace*{1cm}
  \huge{\arrowtitle}

  \vspace*{0.2cm}
  \LARGE{\arrowtype}
  \vspace*{1cm}

  \Large{Service ID: \textit{``\arrowid''}}
  \vspace*{\fill}

  

  \vspace*{1cm}
  \vspace*{\fill}

  % Front Page Abstract
  \begin{abstract}
    This document describes the \texttt{workflowLease} service of the Eclipse Arrowhead workflow support systems, which grants exclusive, expiring use of shared resources such as tools. It is provided by the WorkflowManager system, service version 1.0.0, documentation edition 5.4.0.
  \end{abstract}

  \vspace*{1cm}
\end{center}

\newpage
%%

%% Table of Contents
\tableofcontents
\newpage
%%

\section{Overview}
\label{sec:overview}
The \texttt{workflowLease} service is a FIFO lease manager for shared resources, typically tools
in a tool crib used by several production lines. A step with lease mode \texttt{reserve} reserves the resource; the
returned token is stored in its result and read by the tool-using steps, which \oref{renew} the lease immediately before
each use and fail without a valid lease; a \texttt{release} step releases it. A lease that is not renewed in time
expires, and the next queued reservation is granted, so a failed holder cannot block other sessions forever.

The rest of this document is organized as follows.
In Section~\ref{sec:operations}, we describe the abstract message operations provided by the service.
In Section~\ref{sec:model}, we end the document by presenting the data types used by the mentioned operations.

\subsection{Significant Prior Art}
Proposition~1 of~\cite{varga2026} (exclusive tool use) and the lease variant discussed there and in
the workflow kits. Proposition~1 is proved under the assumption that the tool manager grants one session at a time; for
the lease variant it additionally assumes that every tool step runs inside a lease renewed immediately before it starts
and that the lease is longer than the longest tool step (validator check L04).

\subsection{How This Service Is Meant to Be Used}
The executor running a \texttt{reserve} step calls \oref{reserve}; when the answer is \texttt{QUEUED} it calls
\oref{poll} with the ticket until \texttt{GRANTED} (the call is pending, so the step stays running). Steps in mode
\texttt{use} call \oref{renew} with the token before the invocation. The \texttt{release} step calls \oref{release}.
When a session fails and \texttt{workflow.\allowbreak{}lease.\allowbreak{}release-\allowbreak{}on-\allowbreak{}failure} is set, the WorkflowManager releases its leases
(compensation extension).

\subsection{Important Delimitations and Dependencies}
\label{sec:delimitations}
\begin{itemize}
\item Mutual exclusion is provided by this service, not by the workflow engine. Exclusive use on runs is checked in verification but not proved for the lease variant.
\item Tokens are capabilities: whoever holds a valid token can renew and release. Tokens travel only in $\xi$ through \texttt{workflow\allowbreak{}Context}.
\item Leases are local to one WorkflowManager. An external crib may implement the same contract.
\end{itemize}

\paragraph{Authorization.} Each operation is authorized separately.

\newpage
\section{Service Interface}
\label{sec:operations}
This section describes the interface of the \texttt{workflowLease} service. Each subsection names an abstract operation,
an input type and an output type, in that order. The input type is named inside parentheses, while the output type is
preceded by a colon. Input and output types are only denoted when accepted or returned, respectively, by the operation
in question. All abstract data types named in this section are defined in Section~\ref{sec:model}. The concrete
interface (HTTP/TLS/JSON) is specified in the IDD of the same name.

\begin{figure}[ht!]
  \centering
  \resizebox{0.92\textwidth}{!}{
\begin{tikzpicture}
\node[ahsys] (prov) at (4.6,0) {WorkflowManager\\[-1pt]{\normalfont\scriptsize $\ll$provider$\gg$}};
\node[ahsrv] (srv) at (1.6,0) {workflowLease};
\draw[thick,ArrowheadPurple] (srv.east) -- (prov.west);
\node[ahext] (c0) at (-3.0,0.00) {WorkflowExecutor\\[-1pt]{\scriptsize $\ll$consumer$\gg$}};
\draw[ahflow] (c0.east) -- (srv.west);
\end{tikzpicture}
  }
  \caption{SysML block view of the \texttt{workflowLease} service: provider and typical consumers.}
  \label{fig:SD}
\end{figure}

The following interface operations are available.

{\small
\begin{longtable}{|p{4.4cm}|p{5.6cm}|p{5.8cm}|}

\hline
\rowcolor{gray!33} Operation & Input & Output \\ \hline
\endhead
\oref{reserve} & \mref{LeaseRequest} & \mref{LeaseResponse} \\ \hline
\oref{poll} & \mref{LeaseRequest} & \mref{LeaseResponse} \\ \hline
\oref{renew} & \mref{LeaseRequest} & \mref{LeaseResponse} \\ \hline
\oref{release} & \mref{LeaseRequest} & \mref{LeaseResponse} \\ \hline
\end{longtable}
}

\osubsection{reserve}{LeaseRequest}{LeaseResponse}
Queues a reservation and grants it when the resource is free (FIFO).
\osubsection{poll}{LeaseRequest}{LeaseResponse}
Expires overdue leases, grants the ticket if it is first in line and the resource is free, and returns its state.
\osubsection{renew}{LeaseRequest}{LeaseResponse}
Extends a granted lease by its duration.
\osubsection{release}{LeaseRequest}{LeaseResponse}
Releases a granted lease; the next queued reservation can be granted.

\newpage
\section{Information Model}
\label{sec:model}
Here, all data objects that can be part of the \texttt{workflowLease} service are listed in alphabetic order. Each
subsection, which describes one type of object, begins with the \textit{struct} keyword (a collection of named fields,
each with its own data type), the \textit{enum} keyword (a closed set of values) or the \textit{alias} keyword. As a
complement to the explicitly defined types in this section, there is also a list of implicit primitive types in
Section~\ref{sec:model:primitives}.

\msubsection{struct}{LeaseRequest}
Request of the lease operations. Which fields are used depends on the operation.

{\small
\begin{longtable}{|p{3.9cm}|p{4.4cm}|p{7.5cm}|}

\hline
\rowcolor{gray!33} Field & Type & Description \\ \hline
\endhead
\texttt{resource} & \pref{Name} & \oref{reserve}: the resource. Optional. \\ \hline
\texttt{sessionId} & \pref{SessionId} & \oref{reserve}: the session that will hold the lease (released on failure). Optional. \\ \hline
\texttt{duration\allowbreak{}Seconds} & \pref{Duration} & \oref{reserve}: lease duration; the configured default when absent or not positive. Optional. \\ \hline
\texttt{ticket} & \pref{String} & \oref{poll}: the ticket of a queued reservation. Optional. \\ \hline
\texttt{token} & \pref{String} & \oref{renew} and \oref{release}: the lease token. Optional. \\ \hline
\end{longtable}
}

\msubsection{struct}{LeaseResponse}
State of a reservation.

{\small
\begin{longtable}{|p{3.9cm}|p{4.4cm}|p{7.5cm}|}

\hline
\rowcolor{gray!33} Field & Type & Description \\ \hline
\endhead
\texttt{status} & \mref{LeaseStatus} & State. \\ \hline
\texttt{resource} & \pref{Name} & Resource. Optional. \\ \hline
\texttt{token} & \pref{String} & The lease token while \texttt{GRANTED}; it is the capability to use the resource. Optional. \\ \hline
\texttt{ticket} & \pref{String} & Ticket for \oref{poll}. Optional. \\ \hline
\texttt{leaseUntil} & \pref{DateTime} & Expiry of a granted lease. Optional. \\ \hline
\texttt{message} & \pref{String} & Explanation, e.g. \texttt{No valid lease}. Optional. \\ \hline
\end{longtable}
}

\msubsection{enum}{LeaseStatus}
State of a reservation.

{\small
\begin{longtable}{|p{3.9cm}|p{11.9cm}|}

\hline
\rowcolor{gray!33} Value & Meaning \\ \hline
\endhead
\texttt{GRANTED} & The holder may use the resource until \texttt{leaseUntil}. \\ \hline
\texttt{QUEUED} & Waiting; granted in FIFO order. \\ \hline
\texttt{RELEASED} & Released by the holder, or by compensation when the session failed. \\ \hline
\texttt{EXPIRED} & Not renewed in time; the next queued reservation is granted. \\ \hline
\texttt{UNKNOWN} & Unknown ticket or token. \\ \hline
\end{longtable}
}


\subsection{Primitives}
\label{sec:model:primitives}
Types and structures mentioned throughout this document that are assumed to be available to implementations of this
service. The concrete interpretation of each of these types must be provided by any IDD document claiming to implement
this service.

{\small
\begin{longtable}{|p{3.2cm}|p{12.6cm}|}

\hline
\rowcolor{gray!33} Type & Description \\ \hline
\endhead
\pdef{DateTime} & Pinpoints a moment in time (UTC). \\ \hline
\pdef{Duration} & A length of time in seconds. \\ \hline
\pdef{Name} & An identifier of a recipe, action, step, parameter or resource: a letter followed by at most 62 letters, digits, \texttt{-} or \texttt{\_}. \\ \hline
\pdef{SessionId} & The identifier the WorkflowManager gives a session. \\ \hline
\pdef{String} & A chain of characters. \\ \hline
\end{longtable}
}


\newpage
\bibliographystyle{IEEEtran}
\bibliography{bibliography}

\newpage

\section{Revision History}
\subsection{Amendments}

\noindent\begin{tabularx}{\textwidth}{| p{1cm} | p{2.2cm} | p{1.6cm} | X | p{3.6cm} |} \hline
\rowcolor{gray!33} No. & Date & Version & Subject of Amendments & Author \\ \hline
1 & 2026-11-23 & \arrowversion & Initial edition: service description of workflowLease, written from the Java reference implementation (ah5-Workflow-java). & Documentation maintenance \\ \hline
\end{tabularx}

\subsection{Quality Assurance}

\noindent\begin{tabularx}{\textwidth}{| p{1cm} | p{2.2cm} | p{1.6cm} | X |} \hline
\rowcolor{gray!33} No. & Date & Version & Approved by \\ \hline
1 & & \arrowversion & Pending review \\ \hline
\end{tabularx}

\end{document}
